\section{Motivation}
Climate change is one of the most pressing challenges of our time,
and the energy sector is a major contributor to global greenhouse gas emissions.
Renewable resources like wind and solar have been widely adopted, but another renewable resource, the ocean, remains largely untapped.
There exist a wide variety of devices that harness energy from ocean waves, known as wave energy converters (WECs).
However, it remains unknown what size, shape, operation, and design tradeoffs of a WEC will most effectively contribute to global decarbonization efforts.
This requires the examination and strategic improvement of both the technical and economic performance of wave energy devices.
Several attempts have been made to optimize WECs, but they have been limited in scope, rigor, and relevance.

This thesis aims to perform a wave energy design optimization study that is genuinely meaningful in its outcomes.
In this context, ``meaningful'' is defined by 
\ifdefined\DECIDER
four
\else
three
\fi
criteria.
First, the analysis must be \textbf{comprehensive}, meaning that it spans a sufficiently expansive design space and captures the relevant multidisciplinary couplings inherent in wave energy systems.
Second, it must be \textbf{methodologically rigorous}, meaning that we understand the mathematical structure and numerical behavior of the models sufficiently well to ensure that the optimization results are accurate, convergent, and reproducible.
Third, it must be \textbf{decision-relevant}, in the sense that it yields actionable insights that inform design direction, rather than merely demonstrating a sophisticated numerical method or producing isolated computational results.
\ifdefined\DECIDER
  Fourth, it must be \textbf{aligned with climate objectives}, such that the optimization targets metrics that are directly relevant to economic feasibility and environmental impact.
\fi

Achieving comprehensiveness requires the development of a modeling and optimization framework that is both highly \textbf{integrated} and computationally \textbf{efficient}.
Accordingly, \Cref{ch:modeling} of this dissertation focuses on the development, validation, and benchmarking of such a framework, which is referred to as MDOcean. 
\Cref{ch:hydrodynamic} then provides a detailed theoretical and computational evaluation of the hydrodynamics model, increasing its rigor and extending its capabilities.
The resulting model is packaged in a software called OpenFLASH.
To obtain decision-relevant insights, the MDOcean framework is then applied in a design optimization study, which constitutes \Cref{ch:optimization} of the dissertation. 
\ifdefined\DECIDER
  Finally, to ensure relevance to climate goals, in \Cref{ch:grid} the framework is extended to incorporate performance metrics that better reflect economic and environmental viability.
\fi
%%%%%%%%%%%%%%%%%%%%%%%%%%%%%%%%%%%%%%%%%%%%%%%%%%%%%%%%%%%%%%%%%%

\section{Contributions}
Contribution 1: System Modeling (see \cref{ch:modeling})

Key message: The semi-analytical nature of the MDOcean modeling framework allows it to span multifarious disciplines, reduce computational cost, and build design intuition while maintaining accuracy.
\begin{itemize}
\subimport{../applied-ocean-research-model/sections}{aor-contributions}
\end{itemize}
%
\noindent
Contribution 2: Hydrodynamic Modeling (see \cref{ch:hydrodynamic})

Key message: Consistent with its mathematical structure, the matched eigenfunction expansion method generalizes to complex axisymmetric geometries and converges quickly, accurately, and predictably.

\begin{itemize}
\item Presents a unifying matched eigenfunction expansion method framework for modelling an arbitrary number of fixed or heaving surface-piercing annular cylinders.
\item Documents the block-bidiagonal sparse matrix structure and its computational implications.
\item Conducts numerical experiments to analyze relationships between spatial discretization, series truncation order, nondimensional geometry, accuracy, and runtime.
\item Systematically explores convergence behaviour and outlines a truncation order prediction procedure achieving within 2\% of the target accuracy.
\item Demonstrates <5\% error for slanted geometries, even for angles as steep as 15◦ from vertical.
\item Achieves 2\% convergence of hydrodynamic coefficients 10x faster using 100x smaller matrices vs. standard boundary element method solver.
\end{itemize}
%
\noindent
Contribution 3: Design Optimization (see \cref{ch:optimization})

Key message: Multidisciplinary design optimization exploits careful problem formulation and full coupling to reveal the way that WEC shape and size mediate the tradeoff between cost and power.
\begin{itemize}
\subimport{../renewable-energy-mdo/sections}{re-contributions}
\end{itemize}
%
\ifdefined\DECIDER
  \noindent 
  Contribution 4: Grid Impact Synthesis (see \cref{ch:grid-value})

  Key message: Coupling design optimization with capacity expansion modeling and life cycle assessment can guide early-stage technologies to fulfill system decarbonization goals.
  \begin{itemize}
  \item Identifies features that matter versus don't matter for the integration of grid value and climate change objectives into design optimization: avoided capital cost matters, avoided operational cost matters a bit but less, seasonal variation is more important than hourly variation

  \item Argues that wave energy can provide grid planners/operators with substantial value and that designing for grid value can bring WECs closer to economic viability, though this can only occur if the energy market design appropriately monetizes and incentivizes this

  \item Articulates the connection between energy grid planning optimization (capacity expansion, economic dispatch, and unit commitment problems) and technology specific design optimization and control co-design

  \item Proposes three methods to structure the above optimization (CEM in the loop, CEM sweep with reduced-order WEC design, and price-taker assumption with value matrix) and compares their relevance, fidelity, and limitations

  \item Expands on the implementation of CEM sweep with reduced-order WEC design, shows that a second-order WEC model is insufficient for the above purposes, and proposes other reduced-order models to be explored in future work

  \item Expands on the implementation of the price-taker assumption with value matrix methodology: derive how capacity expansion model outputs can be used to obtain a value matrix which intuitively shows the variation of grid effects with sea state

  \item Shows that under the price taker assumption, the convexity and structure of the net value of energy problem is identical to that of the levelized cost of energy problem, just the weight of power in each sea state changes

  \item Identifies steel, fiberglass, and maintenance fuel as the environmental drivers that can be optimized at the early design stage.
  Shows that the environmental objectives have the same structure as the economic ones, just again with different weights

  \end{itemize}
\fi

\section{List of Research Outputs}\label{sec:research-outputs}
The following publications, presentations, and software have resulted from the work in this dissertation.
$J$ denotes scholarly journal articles, $C$ denotes peer-reviewed conference papers, $P$ denotes presentations, $R$ denotes unpublished technical reports, $O$ denotes open-source code repositories, and $F$ denotes funding awarded from grants or technical competitions.

\newcounter{jcount}
\newcounter{ccount}
\newcounter{pcount}
\newcounter{rcount}
\newcounter{ocount}
\newcounter{fcount}
\newenvironment{typedlist}
{\begin{list}{}{\leftmargin=2em}}
{\end{list}}

\newcommand{\titem}[1]{%
  \ifx#1r
    \stepcounter{rcount}%
    \item[\textbf{R\arabic{rcount}.}]
  \else\ifx#1j
    \stepcounter{jcount}%
    \item[\textbf{J\arabic{jcount}.}]
  \else\ifx#1c
    \stepcounter{ccount}%
    \item[\textbf{C\arabic{ccount}.}]
  \else\ifx#1p
    \stepcounter{pcount}%
    \item[\textbf{P\arabic{pcount}.}]
  \else\ifx#1o
    \stepcounter{ocount}%
    \item[\textbf{O\arabic{ocount}.}]
  \else\ifx#1f
    \stepcounter{fcount}%
    \item[\textbf{F\arabic{fcount}.}]
  \fi\fi\fi\fi\fi
}
%
\noindent
Unified Dissertation Framework:
\begin{typedlist}
    \titem{p} McCabe, ``Leveraging Semi-Analytical Modeling, Multidisciplinary Design Optimization, and System Value Metrics to Advance Wave Energy Converter Viability,'' dissertation defense, May 2026.
    \titem{o} McCabe, Dietrich, Long, Ren, Murphy, and Haji. ``MDOcean.'' Zenodo, May 7, 2026. \url{https://doi.org/10.5281/zenodo.13997243}, \url{https://github.com/symbiotic-engineering/MDOcean}.
    \titem{p} McCabe, ``System Optimization of Ocean Wave Energy Converters,'' Cornell Systems Engineering Ezra Emerging Scholar Seminar, December 2025.
    \titem{f} McCabe, National Science Foundation Graduate Research Fellowship, 2021-2026.
\end{typedlist}
%
\noindent
Contribution 1: System Modeling
\begin{typedlist}
  \titem{j} McCabe, Dietrich, and Haji, ``Development, Validation, and Benchmarking of a Multidisciplinary Semi-Analytical Model for Wave Energy Converters,'' in preparation to submit to Applied Ocean Research in June 2026. \url{https://calkit.io/symbiotic-engineering/mdocean/publications#pubs/applied-ocean-research-model/main.pdf}.
  \titem{c} McCabe and Haji, 2024. ``Force-Limited Control of Wave Energy Converters using a Describing Function Linearization.'' Int. Federation of Automatic Control: Control Applications in Marine Systems, Robotics, \& Vehicles. \url{https://doi.org/10.1016/j.ifacol.2024.10.093}.
\end{typedlist}
%
\noindent
Contribution 2: Hydrodynamic Modeling
\begin{typedlist}
  \titem{j} Bimali, McCabe, Treacy, Khanal, Lo, and Haji. ``Matrix structure and convergence behaviour of the matched eigenfunction method for computing heave wave forces on generalized concentric bodies.'' In review for the Journal of Fluid Mechanics, submission JFM-2026-1453. \url{https://doi.org/10.48550/arXiv.2605.19730}.
  \titem{j} Best, Khanal, McCabe, Jiang, Treacy, and Haji, ``OpenFLASH: An open-source flexible library for analytical and semi-analytical hydrodynamics calculations.'' In review for the Journal of Open Source Software, submission 10.21105.joss.09473. \url{https://joss.theoj.org/papers/842d9827ef856fd88af7128f707f265b}.
  \titem{o} Best, Hope, Rebecca McCabe, Kapil Khanal, Yinghui Bimali, En Lo, Ruiyang Jiang, John Fernandez, Collin Treacy, and Maha N. Haji. ``OpenFLASH.'' Zenodo, June 6, 2026. \url{https://doi.org/10.5281/zenodo.17453418}, \url{https://github.com/symbiotic-engineering/OpenFLASH}.
  \titem{c} McCabe, Khanal, and Haji, 2024. ``Open-Source Toolbox for Semi-Analytical Hydrodynamic Coefficients via the Matched Eigenfunction Expansion Method,'' University Marine Energy Research Community Conference. Award for best communication of complex mathematics. \url{https://doi.org/10.5281/zenodo.14504016}.
  \titem{p} Khanal, McCabe, and Haji, ``Gradient-based Design Optimization of Concentric Cylindrical Offshore Structures,'' SIAM Conference on Optimization, June 2023.
\end{typedlist}
%
\noindent
Contribution 3: Design Optimization
\begin{typedlist}
    \titem{j} McCabe, Dietrich, and Haji, ``Leveraging Multidisciplinary Design Optimization to Advance Wave Energy Converter Viability,'' in preparation to submit to Renewable Energy in June 2026. \url{https://calkit.io/symbiotic-engineering/mdocean/publications#pubs/renewable-energy-mdo/mdocean.pdf}.
    \titem{p} McCabe, Dietrich, and Haji, ``MDOcean: an Open-Source Software for Efficient Simulation and Multidisciplinary Design Optimization of WECs,'' University Marine Energy Research Community Conference, August 2025. Best presentation award. \url{https://doi.org/10.5281/zenodo.20546351}.
    \titem{c} McCabe, Murphy, and Haji, 2022, ``Multidisciplinary Optimization to Reduce Cost and Power Variation of a Wave Energy Converter,'' ASME IDETC-CIE: Design Automation Conference. \url{https://doi.org/10.1115/DETC2022-90227}.
\end{typedlist}
%
\newcommand{\contributionsDecider}{%
\begin{typedlist}
    \titem{c} McCabe, Dietrich, Yang, Long, Kuan, Buccino, Liu, and Haji, 2025. ``WEC Optimization to Maximize Grid Economic Value and Avoided Emissions,'' University Marine Energy Research Community Conference. \url{https://doi.org/10.5281/zenodo.20546181}.
    \titem{c} McCabe, Dietrich, Liu, and Haji, 2023. ``System Level Techno-Economic and Environmental Design Optimization for Ocean Wave Energy.'' ASME IDETC-CIE: Design Automation Conference. \url{https://doi.org/10.1115/DETC2023-114607}.
    \titem{p} McCabe and Haji, ``Value Metrics and Global Impact Potential of Wave Energy,'' University Marine Energy Research Community Conference, September 2022.
\end{typedlist}
}
\ifdefined\DECIDER
  \noindent
  Contribution 4: Grid Impact Synthesis
  \contributionsDecider
\fi

The author has also contributed to the following collaborative publications and projects that are not included in the core contributions of the dissertation.

\ifdefined\DECIDER
\else
  \noindent
  Grid and Environmental Impact Synthesis:
  \contributionsDecider
\fi

\noindent
Wave-Powered Offshore Aquaculture:
\begin{typedlist}
    \titem{c} Hasankhani, McCabe, Ewig, Won, and Haji, 2023. ``Conceptual Design and Optimization of a Wave-Powered Offshore Aquaculture Farm,'' International Ocean and Polar Engineering Conference. \url{https://onepetro.org/ISOPEIOPEC/proceedings/ISOPE23/All-ISOPE23/ISOPE-I-23-112/524501}.
    \titem{c} Hasankhani, Ewig, McCabe, Won, and Haji, 2023. ``Marine Spatial Planning of a Wave-Powered Offshore Aquaculture Farm in the Northeast U.S.'' IEEE Oceans – Limerick. \url{https://doi.org/10.1109/OCEANSLimerick52467.2023.10244332}.
\end{typedlist}
%
\noindent
Wave-Powered Carbon Sequestration:
\begin{typedlist}
    \titem{p} Srivastava, Pokigo, Patel, Olson, Goswami, Doherty, Weng, McCabe, DeGoede, and Haji, ``Wave-Driven Carbon Sequestration: System Requirements and Modeling,'' University Marine Energy Research Community Conference, August 2025. \url{https://doi.org/10.5281/zenodo.20533302}.
    \titem{p} DeGeode, McCabe, Vitale, Pokigo, Srivastava, Edwards, and Haji, ``Offshore Carbon Sequestration Using Direct Mechanical Wave Energy,'' University Marine Energy Research Community Conference, August 2024. \url{https://doi.org/10.5281/zenodo.15168281}.
    \titem{f} U.S. Department of Energy, Power at Sea Prize.
\end{typedlist}
%
\noindent
Wave Energy Prototyping:
\begin{typedlist}
  \titem{j} Vitale, McCabe, Brundan, Mandalam, Alonso Munera, and Haji, 2026. ``Design, build, and analysis of small-scale wave energy converter prototypes.'' ASME Journal of Mechanical Design. \url{http://doi.org/10.1115/1.4070757}.
  \titem{p} Vitale, Brundan, Mandalam, McCabe, and Haji, ``Experimental Investigation of Heterogeneous Wave Energy Converter Arrays: Open-Source Data and Performance Comparisons,'' University Marine Energy Research Community Conference, August 2025. \url{https://doi.org/10.5281/zenodo.20516358}.
\end{typedlist}